\documentclass[10pt,a4paper]{amsart}
\usepackage[margin=19mm]{geometry}
\usepackage{amsmath,amssymb,mathtools}
\usepackage[scr=rsfs,cal=euler]{mathalfa}
\usepackage{xfrac}
\usepackage[unicode,hidelinks,bookmarks=false]{hyperref}
\newcommand{\ie}{i.e.\ }
\newcommand{\cf}{cf.\ }
\newcommand{\resp}{respectively\ }
\newcommand{\Subsection}{\subsection}
\providecommand{\nobreakdash}{\nobreak\hbox{-}}
\setlength{\emergencystretch}{2em}
\multlinegap0pt
\usepackage{bm}
\usepackage{bbm}
\usepackage{dsfont}
\usepackage{xspace}
\usepackage{cancel}
\usepackage{mathtools}
\usepackage{amsthm}
\makeatletter
\@ifundefined{theorem}{\newtheorem{theorem}{Theorem}}{}
\@ifundefined{proposition}{\newtheorem{proposition}[theorem]{Proposition}}{}
\makeatother
\DeclareMathOperator*{\argmax}{arg\,max}
\newcommand{\E}{\mathbb{E}}
\newcommand{\R}{\mathbbm{R}}
\newcommand{\hattheta}{\hat{\theta}}
\newcommand{\mcS}{\mathcal{S}}
\renewcommand{\P}{\mathbb{P}}
\title[Risk reversal for least squares estimators under nested conv]{Risk reversal for least squares estimators under nested convex constraints}
\author{Omar Al-Ghattas}
\date{Source: arXiv:2601.16041v2 (CC BY 4.0)}
\begin{document}
\maketitle

\noindent\emph{Source and licence.} Risk reversal for least squares estimators under nested convex constraints. Version arXiv:2601.16041v2 (CC BY 4.0), distributed under the Creative Commons Attribution 4.0 International licence, as recorded in the arXiv licence metadata for this exact version and shown on the abstract page https://arxiv.org/abs/2601.16041v2. Full rights notice: licenses/arxiv-2601.16041-CC-BY-4.0.txt.\par\smallskip

\noindent\textit{Editorial scope.} Counted unit: the Proposition whose source label is \texttt{prop:quant-diverging-noise-polytope} in the article (source0.tex lines 4266--4286, source label \texttt{prop:quant-diverging-noise-polytope}), delivered together with the article's own complete original proof of it (source0.tex lines 4288--4497, one \texttt{proof} environment of 210 lines). No mathematics is written, repaired or re-derived: statement and proof are the article's own text line for line. Required context retained uncounted: thm:large-sigma-projection (lines 1709--1727). Every definition, display and label of those lines is retained unchanged. Omitted from this package and not redistributed: 1--1708, 1728--4265, 4287--4287, 4498--4661. Nothing inside a retained excerpt is omitted or altered: every retained body is delivered exactly as the article's lines are written, so no omission receipt of any kind accompanies this package. Citations: the retained excerpts carry no citation command at all, so no bibliography entry of the article is transcribed and no reference is redistributed. Reference adaptation: none was needed and none was made; every reference inside the delivered unit resolves inside it, so no unresolved reference or citation is produced. Printed slips kept literal: no printed word, symbol, hypothesis or number of the counted statement or of its proof was altered. Layout and class: the article's own class is \texttt{article} with packages including amsthm, amsmath, amsfonts, amssymb, natbib, hyperref, graphicx, geometry, fontenc, bm, xcolor, array, mathtools, caption; the wrapper is the lane's pinned \texttt{amsart} editorial wrapper at 10pt on A4, which restates the theorem-like environments the retained text needs and adds the article's own macro definitions that the retained text needs (\texttt{\textbackslash E}, \texttt{\textbackslash P}, \texttt{\textbackslash R}, \texttt{\textbackslash argmax}, \texttt{\textbackslash hattheta}, \texttt{\textbackslash mcS}), with extra wrapper packages (bm, bbm, dsfont, xspace, cancel, mathtools). The delivered numbering is the unit's own. Assets: no retained span contains a figure environment, a graphics-inclusion command, a PDF-page inclusion, a table or a TikZ picture; no image file is copied into this package, the package declares no asset dependency and no image file is redistributed with it. Licence: version arXiv:2601.16041v2 of this article is distributed under the Creative Commons Attribution 4.0 International licence, as recorded in the arXiv licence metadata for this exact version and shown on the abstract page https://arxiv.org/abs/2601.16041v2. A full rights notice is delivered at licenses/arxiv-2601.16041-CC-BY-4.0.txt. Characters: every retained excerpt is pure ASCII, so the delivered engine, XeLaTeX with its default Latin Modern text font, reports no missing character.
\begin{theorem}\label{thm:large-sigma-projection}
Let $\Theta\subseteq\R^d$ be nonempty, compact and convex and fix $\theta^\star\in\Theta$.
Let $Z\sim N(0,I_d)$ and for $\sigma>0$ define $\hattheta_\sigma := \Pi_\Theta(\theta^\star+\sigma Z)$. Set
\[
U:=Z/\|Z\|,
\qquad
F_\Theta(u):=\argmax_{\theta\in\Theta}\langle \theta,u\rangle.
\]
Equivalently, $F_\Theta(u)$ is the exposed face of $\Theta$ in direction $u$. 
Then, almost surely as $\sigma\to\infty$,
\begin{align}\label{eq:large-sigma-a.s}
\hattheta_\sigma \to \Pi_{F_\Theta(U)}(\theta^\star),
\end{align}
and consequently, for $R_\infty(\theta^\star; \Theta) := \E\|\Pi_{F_\Theta(U)}(\theta^\star)-\theta^\star\|^2$,
\begin{align}
R_\sigma(\theta^\star; \Theta) \to R_\infty(\theta^\star; \Theta).
\label{eq:large-sigma-risk}
\end{align}
\end{theorem}

\begin{proposition}\label{prop:quant-diverging-noise-polytope}
Fix $d\ge2$, and for $m \ge 2$ let $v_1,\ldots,v_m\in\R^d$ be the distinct vertices of a compact, convex polytope $\Theta=\operatorname{conv}(v_1,\ldots,v_m)\subseteq\R^d$. Set
\[
    D_{\Theta}:=\operatorname{diam}(\Theta),
    \qquad
    \delta_\Theta:=\min_{i\neq j}\|v_i-v_j\|.
\]
Then there exists a constant $C_d<\infty$, depending only on $d$, such that, uniformly over $\theta^\star\in\Theta$,
\[
0
\le
R_\infty(\theta^\star;\Theta)-R_\sigma(\theta^\star;\Theta)
\le
\frac{
    C_d \binom{m}{2}D_\Theta^4
}{
    \delta_\Theta\sigma
},
\qquad \sigma>0.
\]
\end{proposition}

\begin{proof}
Fix $\theta^\star\in\Theta$. Write $Z=\|Z\|U$, where $U\sim\operatorname{Unif}(\mcS^{d-1})$ and $U$ is independent of $\|Z\|$. For $u\in\mcS^{d-1}$, set 
\[
F_\Theta(u) := \operatorname*{arg\,max}_{\theta\in\Theta} \langle u,\theta\rangle,
\qquad
Q(u):=\Pi_{F_\Theta(u)}(\theta^\star).
\] 
Then Theorem~\ref{thm:large-sigma-projection} gives $R_\infty(\theta^\star;\Theta)
=
\E\|Q(U)-\theta^\star\|^2$.
On the other hand, writing $P_\tau(u):=\Pi_\Theta(\theta^\star+\tau u)$ for a given radius $\tau>0$, the finite-noise risk can be written as 
\[
R_\sigma(\theta^\star;\Theta) 
= \E \bigl\| P_{\sigma\|Z\|}(U)-\theta^\star \bigr\|^2. 
\]

We now prove two deterministic bounds relating the finite-radius projection $P_\tau(u)$ and the limiting exposed-face projection $Q(u)$ for any $u \in \mcS^{d-1}$. 
First, note that $P_\tau(u)$ maximizes
\[
    \theta\mapsto
    \langle u,\theta\rangle
    -
    \frac{1}{2\tau}\|\theta-\theta^\star\|^2
\]
over $\Theta$. Comparing this objective at $P_\tau(u)$ and at
$Q(u)\in F_\Theta(u)$, we obtain
\[
    h_\Theta(u)-\langle u,P_\tau(u)\rangle
    \le
    \frac{
    \|Q(u)-\theta^\star\|^2
    -
    \|P_\tau(u)-\theta^\star\|^2
    }{2\tau}.
\]
Since the left-hand side is nonnegative, we have
\[
    0
    \le
    \|Q(u)-\theta^\star\|^2
    -
    \|P_\tau(u)-\theta^\star\|^2.
\]
Second, we argue that the two projections coincide unless $u$ lies close to a
boundary where the exposed face $F_\Theta(u)$ changes.

To make this precise, let $V(\Theta)$ denote the vertex set of $\Theta$, and define 
\[
\gamma_\Theta(u)
:=
h_\Theta(u)
-
\max\{ \langle u,v\rangle: v\in V(\Theta),\ v\notin F_\Theta(u) \},
\qquad
h_\Theta(u):=\max_{\theta\in\Theta}\langle u,\theta\rangle,
\] 
with the convention that $\gamma_\Theta(u)=+\infty$ if all vertices lie in $F_\Theta(u)$. We call $\gamma_\Theta(u)$ the directional margin. It is the gap between the support value $h_\Theta(u)$ and the largest value of
$\langle u,v\rangle$ among vertices not belonging to the exposed face
$F_\Theta(u)$. We claim that $P_\tau(u)$ and $Q(u)$ differ only on the small-margin set
\[
\left\{ u\in\mcS^{d-1}: \gamma_\Theta(u)\le \frac{D_{\Theta}^2}{\tau} \right\}. 
\]

Let $y=\theta^\star+\tau u$. Since $P_\tau(u)=\Pi_\Theta(y)$, to prove $P_\tau(u)=Q(u)$ it is enough
to show that $Q(u)$ satisfies the variational characterization of the
projection of $y$ onto $\Theta$. Since $Q(u)\in F_\Theta(u)\subseteq\Theta$,
this amounts to verifying that
\[
    \langle y-Q(u),\theta-Q(u)\rangle\le0
    \qquad
    \text{for all }\theta\in\Theta.
\]
Write $\theta=\sum_{i=1}^m\lambda_i v_i$, where
$\lambda_i\ge0$ and $\sum_i\lambda_i=1$. Let $I:=\{i:v_i\in F_\Theta(u)\}$. For $i\in I$, since $v_i\in F_\Theta(u)$ and
$Q(u)=\Pi_{F_\Theta(u)}(\theta^\star)$,
\[
    \langle \theta^\star-Q(u),v_i-Q(u)\rangle\le0,
    \qquad
    \langle u,v_i-Q(u)\rangle=0.
\]
Here the first inequality is the projection optimality condition, while the
second identity follows because both $v_i$ and $Q(u)$ lie in the exposed
face $F_\Theta(u)$. 
For $i\notin I$, by Cauchy-Schwarz,
\[
    \langle \theta^\star-Q(u),v_i-Q(u)\rangle
    \le
    \|\theta^\star-Q(u)\|\,\|v_i-Q(u)\|
    \le
    D_{\Theta}^2,
\]
while, by the definition of $\gamma_\Theta(u)$,
\[
    \langle u,v_i-Q(u)\rangle
    =
    \langle u,v_i\rangle-h_\Theta(u)
    \le
    -\gamma_\Theta(u).
\]
Therefore
\[
    \langle y-Q(u),\theta-Q(u)\rangle
    =
    \sum_{i=1}^m
    \lambda_i
    \langle \theta^\star-Q(u)+\tau u,v_i-Q(u)\rangle
    \le
    \sum_{i\notin I}
    \lambda_i
    \bigl(D_{\Theta}^2-\tau\gamma_\Theta(u)\bigr)
    \le0
\]
whenever $\tau\gamma_\Theta(u)\ge D_{\Theta}^2$. This proves the claim.

Combining the preceding observations gives, for every $u\in\mcS^{d-1}$
and $\tau>0$,
\[
    0
    \le
    \|Q(u)-\theta^\star\|^2
    -
    \|P_\tau(u)-\theta^\star\|^2
    \le
    D_{\Theta}^2
    \mathbf 1
    \left\{
        \gamma_\Theta(u)
        \le
        \frac{D_{\Theta}^2}{\tau}
    \right\}.
\]

It remains to control the probability of small margins. If
$\gamma_\Theta(u)\le s$, then there exist distinct vertices $v_i,v_j$ such
that $|\langle u,v_i-v_j\rangle|\le s$. Indeed, one may take $v_i\in F_\Theta(u)$ and a vertex $v_j\notin F_\Theta(u)$
whose value of $\langle u,v_j\rangle$ is within $s$ of the maximum. Thus
\[
    \{u:\gamma_\Theta(u)\le s\}
    \subseteq
    \bigcup_{i<j}
    \{
        u\in\mcS^{d-1}:
        |\langle u,v_i-v_j\rangle|\le s
    \}.
\]
Using rotational invariance, for every nonzero $a\in\R^d$, $\langle U,a\rangle
\stackrel{d}{=}\|a\|U_1$,
where $U_1$ is the first coordinate of $U$. Since $U_1$ has bounded density in a neighborhood of zero, there exists a constant $C_d<\infty$, depending only on $d$, such that for every $t>0$,
\[
    \P(|U_1|\le t)\le C_d\min\{1,t\}.
\]
Therefore, for every $s>0$,
\[
    \P(|\langle U,a\rangle|\le s)
    =
    \P\left(|U_1|\le \frac{s}{\|a\|}\right)
    \le
    C_d\min\left\{1,\frac{s}{\|a\|}\right\}.
\]
Using $\|v_i-v_j\| \ge \delta_{\Theta}$, a union bound then gives
\[
    \P\bigl(\gamma_\Theta(U)\le s\bigr)
    \le
    C_d \binom{m}{2}\frac{s}{\delta_\Theta},
    \qquad s>0.
\]

We now apply the deterministic estimate with
$\tau=\sigma\|Z\|$. Conditioning on $\|Z\|$, we get
\[
\begin{aligned}
0
&\le
R_\infty(\theta^\star;\Theta)
-
R_\sigma(\theta^\star;\Theta)  \\
&\le
D_{\Theta}^2
\E\left[
\P\left(
    \gamma_\Theta(U)
    \le
    \frac{D_{\Theta}^2}{\sigma\|Z\|}
    \,\middle|\, \|Z\|
\right)
\right] \\
&\le
D_{\Theta}^2
\E\left[
    C_d \binom{m}{2}
    \frac{D_{\Theta}^2}{\delta_\Theta\sigma\|Z\|}
\right] \\
&=
C_d \binom{m}{2}
\frac{D_{\Theta}^4}{\delta_\Theta\sigma}
\E\|Z\|^{-1}.
\end{aligned}
\]
The bound is uniform in $\theta^\star$. Since $d\ge2$,
\[
    \E\|Z\|^{-1}
    =
    2^{-1/2}
    \frac{\Gamma((d-1)/2)}{\Gamma(d/2)}
    <
    \infty.
\]
Absorbing the finite factor $\E\|Z\|^{-1}$, which depends only on $d$, into $C_d$
gives the claimed bound.
\end{proof}

\end{document}
